% PLAN (15-20 pages). Tutorial background first, then the related work of the paper, which is input unchanged at the end.
%   2.1 Chest radiography and bone overlap: image formation, attenuation, why ribs and clavicles hide lesions.
%   2.2 Dual-energy subtraction: principle, hardware, dose, SZCH-X-Rays as a public DES data set (sun2026bsldm, sun2026debonedit).
%   2.3 Learning-based bone suppression: MTANN (suzuki2006mtann), CNNs (yang2017cascade, gusarev2017deep), diffusion (chen2024bsdiff, sun2026bsldm,
%       sun2025gllcm, sun2026debonedit), CT-supervised (angaitkar2026anatomy).
%   2.4 Architectures used here: Swin Transformer V2 (liu2022swinv2), UNet++ (zhou2018unetpp), CBAM (woo2018cbam); optimiser (loshchilov2019adamw).
%   2.5 Image-quality and detection metrics: PSNR, SSIM (wang2004ssim), LPIPS (zhang2018lpips), CNR, FROC and CPM (niemeijer2011cpm), AUROC and
%       DeLong (delong1988auc).
%   2.6 Uncertainty: deep ensembles (lakshminarayanan2017ensembles), split conformal prediction (lei2018conformal, angelopoulos2023conformal).
%   2.7 Statistics: paired bootstrap, cluster bootstrap, Holm correction, non-inferiority, pre-specified analysis plans.
%   2.8 Related work as in the paper (input below).
% Section 2.3 describes each cited paper only as far as its abstract (PubMed or arXiv, read 2026-11-03) supports; check against the PDFs before submission.
This chapter collects the background that the later chapters assume. It is written for a reader who knows machine learning but not the physics
of chest radiography or the conventions of this field. Sections~\ref{sec:bg_cxr} and~\ref{sec:bg_des} explain how a chest radiograph is formed
and why bone overlap is a problem; Section~\ref{sec:bg_learned} surveys learned bone suppression; Section~\ref{sec:bg_arch} introduces the network
components and training techniques used in this study; Sections~\ref{sec:bg_metrics} to~\ref{sec:bg_stats} define the metrics, the uncertainty
estimates and the statistical procedures. The related work of the paper closes the chapter (Section~\ref{sec:related}). No result of this study
appears here.

\section{Chest radiography and bone overlap}
\label{sec:bg_cxr}
\paragraph{Projection imaging.} A chest radiograph is a projection image. An X-ray tube emits a beam that passes through the thorax and is
recorded by a detector on the other side. In the standard posteroanterior (PA) view the patient faces the detector and the beam enters from the
back; the anteroposterior (AP) view, used at the bedside, reverses the direction. The tube produces a polyenergetic spectrum: bremsstrahlung up to
a maximum photon energy set by the tube voltage (kVp), plus characteristic lines of the anode material, shaped by the filtration in the beam. Each
detector pixel records the energy that reaches it along the rays that end in that pixel. Everything the ray crosses, from the skin of the back to
the skin of the chest, contributes to the same pixel; depth information is lost. This superposition is the root of the problem studied in this
thesis.

\paragraph{Attenuation.} For a narrow monoenergetic beam of energy $E$ crossing a homogeneous slab of thickness $t$, the transmitted intensity
follows the Beer--Lambert law, $I=I_0\,e^{-\mu(E)\,t}$, where $\mu$ is the linear attenuation coefficient of the material. Through a
heterogeneous body the exponent becomes a line integral along the ray $L$:
\begin{equation}
\begin{gathered}
I(E)=I_0(E)\,\exp\!\Big(-\int_L \mu(\mathbf{r},E)\,\mathrm{d}r\Big),\\
p(E)=-\ln\frac{I(E)}{I_0(E)}=\int_L \mu(\mathbf{r},E)\,\mathrm{d}r .
\end{gathered}
\label{eq:bg_beerlambert}
\end{equation}
The log-transformed value $p$ is linear in the attenuation along the ray, so overlapping structures \emph{add}. It is convenient to write
$\mu=(\mu/\rho)\,\rho$, the product of the mass attenuation coefficient $\mu/\rho$, which depends only on the elemental composition and the
energy, and the density $\rho$. If the body is described by a few materials $c$ (for example bone, soft tissue and lung) with fixed compositions,
the line integral splits into a sum over materials,
\begin{equation}
\begin{gathered}
p(E)=\sum_c \Big(\frac{\mu}{\rho}\Big)_{\!c}(E)\;a_c,\\
a_c=\int_L \rho_c(\mathbf{r})\,\mathrm{d}r ,
\end{gathered}
\label{eq:bg_materials}
\end{equation}
where $a_c$ is the areal density (g\,cm$^{-2}$) of material $c$ along the ray. Tabulated mass attenuation coefficients for elements and mixtures
are available from NIST XCOM \citep{berger2010xcom}, and reference compositions of body tissues from ICRU Report~44 \citep{icru1989report44}.
Equation~\eqref{eq:bg_materials} is the basis of everything that follows: of dual-energy subtraction (Section~\ref{sec:bg_des}), of the
CT-derived training data of this study, and of the idea that a radiograph can be written as a soft-tissue image plus a bone image.

\paragraph{Polyenergetic beams and beam hardening.} A real beam has a spectrum $\phi(E)$, and an energy-integrating detector weights each photon
roughly by its energy. The recorded log signal is then
\begin{equation}
p=-\ln\frac{\int \phi(E)\,E\,\exp\!\big(-\sum_c (\mu/\rho)_c(E)\,a_c\big)\,\mathrm{d}E}{\int \phi(E)\,E\,\mathrm{d}E},
\label{eq:bg_poly}
\end{equation}
the model used in Chapter~\ref{ch:data}. Because low-energy photons are absorbed preferentially, the mean energy of the beam rises along its path
(\emph{beam hardening}) and $p$ is no longer exactly linear in the $a_c$: the sum of separately projected materials differs slightly from the
projection of the whole body. Chapter~\ref{ch:data} quantifies this additivity gap for the renderer of this study.

\paragraph{Why bone stands out.} At diagnostic energies (roughly 20 to 150~keV) two interactions dominate attenuation in tissue. Photoelectric
absorption rises steeply with atomic number and falls steeply with energy; Compton scattering depends mainly on electron density and changes
slowly with energy \citep{bushberg2020essential}.
Bone mineral contains calcium and phosphorus, whose atomic numbers are much higher than those of the hydrogen, carbon, nitrogen and oxygen that make
up soft tissue, so bone attenuates more than soft tissue at every energy and its attenuation falls faster with energy. The first property makes
ribs and clavicles bright on a radiograph; the second makes dual-energy subtraction possible.

\paragraph{From detector signal to stored image.} Scattered photons reach the detector from directions other than the primary ray and add a
smooth haze that lowers contrast; anti-scatter grids reduce but do not remove it. The number of detected photons is random, which adds quantum
noise. Before an image is stored, digital systems apply vendor-specific processing (logarithmic or other tone curves, look-up tables, multi-scale
contrast enhancement and edge sharpening) \citep{prokop2003processing,schaeferprokop2008digital}. A stored radiograph is therefore
an unknown, mostly monotone function of the line integral in Equation~\eqref{eq:bg_poly}, combined with local filtering. This is why the training
data of this study pass the exact components through random monotone tone curves (Section~\ref{sec:methods}) and why real-image metrics are also
reported after a per-image tone calibration (Section~\ref{sec:bg_metrics}).

\paragraph{Why ribs and clavicles hide lesions.} On a PA radiograph the posterior and anterior ribs, the clavicles and parts of the scapulae
project over the lung fields, and the anterior and posterior ribs cross each other. A pulmonary nodule is a small region of near-water density
inside the air-filled lung; in the log domain it adds a compact signal of a few millimetres to a few centimetres. A rib adds a signal of similar or
larger amplitude, with sharp cortical edges, at the same spatial scale. In a linear projection the nodule over a rib keeps its own signal, but it
sits on a background that is bright, structured and full of edges, and the eye (or a detection algorithm) has to separate the two. Background
structure of this kind, often called anatomical noise, is generally regarded as a stronger limit on nodule detection in chest radiography than
quantum noise \citep{samei1999detection}. Reader studies support the clinical relevance
of bone overlap: in a study in which fifteen radiologists read the radiographs of 368 patients first unaided and then with software that
suppresses ribs and clavicles, filters noise and equalises contrast, reader sensitivity rose from 49.5\,\% to 66.3\,\% while specificity fell from
96.1\,\% to 91.8\,\%, and 74\,\% of the aided detections were of cancers with at least 70\,\% overlap between bone and nodule \citep{freedman2011lung}.

\paragraph{Bone suppression as source separation.} In the log domain a radiograph can be written as $x=s+b$, a soft-tissue image $s$ plus a bone
image $b$ (up to the beam-hardening gap and the tone curve). Bone suppression estimates $b$, or equivalently $s=x-b$, from $x$ alone. This is a
separation problem with one equation and two unknowns per pixel; it can only be solved with prior knowledge of what bone looks like (elongated,
smoothly curved ribs with bright cortical edges, clavicles of known shape and position) and of what soft tissue looks like. The prior is weakest
exactly where it matters: a compact bright nodule that lies on a rib is consistent with both explanations. A suppressor that has learned that
bright compact structures in the rib region belong to bone can move part of the nodule into $\hat b$ and remove it from the soft image. The audit
of Chapter~\ref{ch:audit} measures this lesion absorption directly.

\section{Dual-energy subtraction}
\label{sec:bg_des}
\paragraph{Principle.} Dual-energy subtraction (DES) separates bone from soft tissue by imaging the same anatomy with two different spectra.
Write $p_L$ and $p_H$ for the log images obtained with a low-energy and a high-energy spectrum, and treat the body as two materials, soft tissue
($\mathrm{s}$) and bone ($\mathrm{b}$), with effective mass attenuation coefficients $\mu^{L}_{\cdot}$ and $\mu^{H}_{\cdot}$ for the two spectra.
Equation~\eqref{eq:bg_materials} then gives two equations in two unknowns:
\begin{equation}
\begin{gathered}
p_L=\mu^{L}_{\mathrm{s}}\,a_{\mathrm{s}}+\mu^{L}_{\mathrm{b}}\,a_{\mathrm{b}},\\
p_H=\mu^{H}_{\mathrm{s}}\,a_{\mathrm{s}}+\mu^{H}_{\mathrm{b}}\,a_{\mathrm{b}} .
\end{gathered}
\label{eq:bg_des_system}
\end{equation}
A weighted subtraction of the log images cancels one material:
\begin{align}
p_{\mathrm{soft}} &= p_H-w_{\mathrm{b}}\,p_L=\big(\mu^{H}_{\mathrm{s}}-w_{\mathrm{b}}\,\mu^{L}_{\mathrm{s}}\big)\,a_{\mathrm{s}},
 & w_{\mathrm{b}}&=\mu^{H}_{\mathrm{b}}/\mu^{L}_{\mathrm{b}}, \label{eq:bg_des_soft}\\
p_{\mathrm{bone}} &= p_H-w_{\mathrm{s}}\,p_L=\big(\mu^{H}_{\mathrm{b}}-w_{\mathrm{s}}\,\mu^{L}_{\mathrm{b}}\big)\,a_{\mathrm{b}},
 & w_{\mathrm{s}}&=\mu^{H}_{\mathrm{s}}/\mu^{L}_{\mathrm{s}}. \label{eq:bg_des_bone}
\end{align}
The system is solvable only if the ratios $w_{\mathrm{b}}$ and $w_{\mathrm{s}}$ differ, that is, if bone and soft tissue change their attenuation
differently with energy; the photoelectric effect in bone mineral provides this difference (Section~\ref{sec:bg_cxr}). The soft-tissue image
$p_{\mathrm{soft}}$ shows lung, vessels, heart and soft-tissue lesions without ribs and clavicles; the bone image $p_{\mathrm{bone}}$ shows the
skeleton and anything else that contains calcium.

\paragraph{What DES does and does not do.} The decomposition is a two-material model. A third material with a different energy dependence is
split between the two images according to its own coefficients, and calcium outside the skeleton is, by construction, bone. A calcified nodule
therefore appears in the DES bone image, which is one of the main clinical uses of the technique: a review of DES chest radiography names the
clearer depiction of calcification, which helps to characterise pulmonary nodules, as its main advantage and lists further uses in the hila,
mediastinum, airways, vessels, pleura and chest wall \citep{kuhlman2006des}. A non-calcified nodule, by contrast, belongs to the soft-tissue
image. This distinction defines what a learned suppressor should do with a soft-tissue lesion over a rib: leave it in $s$. In practice the
weights in Equations~\eqref{eq:bg_des_soft} and~\eqref{eq:bg_des_bone} are tuned empirically, because the spectra are polyenergetic and beam
hardening and scatter make the effective coefficients depend on thickness; cancellation is imperfect and residual bone edges can remain. The
subtraction of two images also removes most of the signal while their noise adds, so DES images are noisier than the standard radiograph and are
usually filtered, which changes their texture and tone.

\paragraph{Hardware, dose and motion.} Two acquisition schemes are in use \citep{schaeferprokop2008digital}. In \emph{dual-exposure} DES the tube voltage is switched between two exposures taken in quick succession; the energy separation is
good, but anything that moves between the two exposures (the heart, the diaphragm, the patient) is misregistered and leaves artefacts in both
images \citep{schaeferprokop2008digital}. In \emph{single-exposure} DES a layered detector with a filter between the layers
records the softer part of the spectrum in front and the harder part behind; there is no motion between the images, but the energy separation is
smaller and the images are noisier. Dual-exposure radiographs exist in research collections; for example, one CNN study trained on 504 cases of
two-exposure DES chest radiographs \citep{yang2017cascade}. DES needs dedicated equipment, and the review cited above states that it requires a
somewhat higher radiation dose than conventional radiography \citep{kuhlman2006des}; authors of learned methods cite equipment cost and
additional exposure as reasons to replace DES by software \citep{chen2024bsdiff}. Most chest radiographs, including all large public
collections used in this study, are single-energy images.

\paragraph{Public paired data.} Learning bone suppression from real images requires radiographs paired with DES images of the same exposure, and
such pairs are scarce. The data sets of the earlier learned methods (Section~\ref{sec:bg_learned}) are not described as public in their
abstracts. The SZCH-X-Rays data set was compiled for BS-LDM, whose abstract describes 818 pairs of high-resolution radiographs and soft-tissue
images from a partner hospital \citep{sun2026bsldm}; it was released with DeBoneDiT, whose abstract describes the release as 741 pairs of
radiographs and DES soft-tissue images and as the largest publicly available high-quality paired data set for bone suppression \citep{sun2026debonedit}.
The BS-LDM paper compiled 831 pairs and excluded 13; DeBoneDiT and GL-LCM report 741 pairs split into 592 training, 74 validation and 75
test pairs, with the same exclusion criteria. None of the three papers explains the difference, and only the 75 test pairs are public. The
images come from one partner hospital, acquired with dual-exposure DES on a GE Discovery XR656 at $2021\times2021$ pixels of 0.194~mm.
The released data set carries a non-commercial licence (Chapter~\ref{ch:data}). In this study the SZCH-X-Rays pairs are used only as
evaluation references, never for training. Three properties of DES references matter for that use. They are themselves estimates, with residual
bone, noise and possible motion artefacts; their tone and texture differ from the radiograph they are paired with; and they come from one
institution. Chapter~\ref{ch:benchmark} therefore reports metrics both as computed and after a per-image gain and offset fit, and treats
agreement with DES as one line of evidence among several.

\section{Learning-based bone suppression}
\label{sec:bg_learned}
This section follows the field roughly in chronological order, grouping closely related studies, and describes each method only as far as its
abstract supports. Input resolution, which most abstracts do not state, is left out. Table~\ref{tab:bg_learned} summarises the learned methods; Section~\ref{sec:related} places them
relative to this study.

\begin{table}[htbp]
\centering\small
\renewcommand{\arraystretch}{1.15}
\caption{Learned bone-suppression methods discussed in Section~\ref{sec:bg_learned}, described only as far as their abstracts support. Rows
follow the order of first appearance (preprint date where there is one); the year is that of the bibliography entry; data sizes are those stated
in the abstracts.}
\label{tab:bg_learned}
\begin{tabular}{@{}p{0.07\linewidth}p{0.21\linewidth}p{0.29\linewidth}p{0.30\linewidth}@{}}
\toprule
Year & Method & Model & Training pairs named in the abstract \\
\midrule
2006 & MTANN \citep{suzuki2006mtann} & multiresolution massive-training ANN (non-linear filter) & DES bone images as teaching images (number not stated) \\
2017 & CamsNet \citep{yang2017cascade} & cascade of multi-scale CNNs in the gradient domain & 504 two-exposure DES cases (404 training, 100 test) \\
2017 & \citet{gusarev2017deep} & two end-to-end networks & not stated in the abstract \\
2024 & BS-Diff \citep{chen2024bsdiff} & conditional diffusion model (U-Net) with an autoencoder-based enhancement module & radiographs and soft-tissue images of 120 patients \\
2026 & BS-LDM \citep{sun2026bsldm} & conditional latent diffusion with a VQ-GAN & SZCH-X-Rays (818 pairs); 241 JSRT pairs \\
2025 & GL-LCM \citep{sun2025gllcm} & latent consistency model with global and local paths & SZCH-X-Rays; JSRT \\
2026 & \citet{angaitkar2026anatomy} & networks that predict one CT-projected component & projections of chest CT; no real pairs \\
2026 & DeBoneDiT \citep{sun2026debonedit} & source-conditioned Brownian-bridge diffusion transformer & SZCH-X-Rays release (741 DES pairs) \\
\bottomrule
\end{tabular}
\end{table}

\paragraph{Massive-training neural networks (2006).} \citet{suzuki2006mtann} trained a multiresolution massive-training artificial neural
network (MTANN), a highly non-linear filter applied separately at three resolution levels, to map a radiograph to its DES bone image. The
network outputs a bone-like image, and subtracting it from the radiograph gives a soft-tissue-like image. On a validation set of 118 radiographs
with nodules and an independent test set of 136 digitised screen-film radiographs from 14 institutions, ribs and clavicles were substantially
suppressed while the visibility of nodules and vessels was maintained; the abstract reports no detection measure. This
predict-the-bone-and-subtract design is the one used in this study.

\paragraph{Reader and CAD studies (2011--2014).} The reader study of \citet{freedman2011lung}, summarised in Section~\ref{sec:bg_cxr}, showed
higher sensitivity with suppression software at some cost in specificity. A later study asked whether computer-aided detection (CAD) adds anything once readers
already have bone-suppressed images: with 300 subjects (111 with a solitary nodule) read by eight observers with PA and lateral radiographs and
bone-suppressed images, reading with CAD marks increased the JAFROC figure of merit, although observers dismissed most of the true-positive CAD
candidates they had missed \citep{schalekamp2014cad}. These studies use commercial software whose algorithm their abstracts do not describe.

\paragraph{Contrast of inserted nodules (2013, 2020).} Two studies measured what suppression does to a nodule whose position is known.
\citet{hogeweg2013suppression} suppressed elongated structures (ribs, clavicles, catheters) with a non-learned method that models intensity
profiles across the structure and separates them by blind source separation. In 100 radiographs with nodules extracted from CT and simulated near
the clavicles, the abstract reports that clavicle contrast fell from 2.46 to 0.55 after suppression, while nodule contrast fell from 1.83 to 1.35:
suppression removed most of the clavicle and some of the nodule. \citet{takagi2020dose} imaged an anthropomorphic phantom with a 12~mm simulated
nodule at 66 positions under several tube voltages and doses; in areas not overlapping the ribs, bone suppression significantly increased nodule
contrast and CNR, and bone-suppressed images at 70~kVp were equal to or better than those at 140~kVp. The abstract does not name the
suppression software. Both studies measure contrast of a known nodule; neither audits a learned suppressor against an exact, additive lesion
signal (Section~\ref{sec:related}).

\paragraph{Convolutional networks (2017).} \citet{yang2017cascade} trained a cascade of multi-scale convolutional networks (CamsNet) to map
the gradients of a radiograph to the gradients of its bone image, refined the prediction from coarse to fine scales, fused the predicted bone
gradients in a maximum-a-posteriori framework to reconstruct the bone image, and subtracted it from the radiograph. Training and testing used 504
cases of real two-exposure DES radiographs (404 for training and 100 for testing) with a pixel size of about 0.194~mm, and the trained model was
also applied to radiographs from other machines, including scanned films. A plausible reason for working in the gradient domain is that the bone
signal varies most at the rib edges, so gradients concentrate the learning problem there. \citet{gusarev2017deep} proposed two end-to-end architectures that treat bones as noise to be
removed and reported that the scheme does not introduce artefacts; their abstract does not describe the training data.

\paragraph{Diffusion models on DES pairs (2024--2026).} Diffusion models generate an image by reversing a gradual noising process; a
\emph{conditional} diffusion model generates the soft-tissue image given the radiograph. BS-Diff \citep{chen2024bsdiff} combines a conditional
diffusion U-Net with an enhancement module that includes an autoencoder, trained on paired radiographs and soft-tissue images of 120 patients
from the authors' hospital, and evaluates it with comparisons, ablations and clinical evaluations. BS-LDM \citep{sun2026bsldm} moves the diffusion
into the latent space of a vector-quantised GAN trained with a multi-level hybrid loss, adds offset noise in the forward process and a temporal
adaptive thresholding in the reverse process to preserve low-frequency soft-tissue information, and introduces SZCH-X-Rays; it also uses 241
pairs from the JSRT database \citep{shiraishi2000jsrt} converted to negative images. GL-LCM \citep{sun2025gllcm} replaces diffusion sampling by a
latent consistency model for speed and combines a global path with a local path guided by lung segmentation, adding a training-free local
guidance to reduce boundary artefacts; it is evaluated on SZCH-X-Rays and JSRT. DeBoneDiT \citep{sun2026debonedit} formulates suppression as a
source-conditioned Brownian-bridge diffusion that starts from the radiograph instead of from noise, uses a transformer denoiser in the latent
space of an autoencoder trained with a feature loss from a pretrained depth-estimation encoder, and releases SZCH-X-Rays. These methods produce
the soft-tissue image directly and are trained on real DES pairs from one institution; generative models of this kind can also produce content
that is not in the input (Section~\ref{sec:related}).

\paragraph{CT-supervised suppression (2026).} \citet{angaitkar2026anatomy} replace real pairs by digitally reconstructed radiographs: chest CT
volumes are decomposed into bone, non-lung soft tissue and lung, each component is projected separately, and the weighted components sum exactly
to the full projection. Models learn to predict one component (bone or lung) and recover the rest by subtraction, and they transfer to real
radiographs without real paired training data. The abstract reports detection gains concentrated on abnormalities with substantial bone overlap and
releases models and projections. Supervision of this kind is exact and scalable, because every CT volume yields a pair, but it carries a domain gap
between projections and real radiographs, and the lesions of the training CTs are whatever the CT collection happens to contain. This study builds
on these projections (Chapter~\ref{ch:v2}).

\paragraph{Summary.} Three sources of supervision appear in the literature: real DES pairs, which are realistic but scarce, noisy and
institution-specific; non-learned models of the bone, which need no pairs but have hand-set assumptions; and CT projections, which are exact and
plentiful but synthetic. Evaluation has used similarity to DES images, reader and CAD studies, downstream detection and the contrast of known
nodules. What a learned suppressor does to a soft-tissue lesion that lies on a rib, measured against an exactly known lesion signal, is the
question of this thesis.

\section{Architectures and optimisation}
\label{sec:bg_arch}
The suppressor of this study is an encoder--decoder network that maps a radiograph $x$ to a bone image $\hat b$; the soft image is
$\hat s=x-\hat b$. Predicting the bone image rather than the soft image keeps the output sparse (zero outside bone) and makes the identity
$x=\hat s+\hat b$ hold by construction. The encoder is a pretrained image backbone that produces a feature pyramid; the decoder merges the pyramid
back to full resolution. The confirmatory dose-response and the ablation use a Swin Transformer V2 Tiny encoder (SwinV2-T); the final model uses a
ConvNeXt V2 Tiny encoder (ConvNeXt-V2-T), chosen by the pre-declared rule described in Section~\ref{sec:experiments}. Both produce four feature maps
at strides 4, 8, 16 and 32, so the same decoder, after a $1\times1$ projection of each map, serves both.

\paragraph{Swin Transformer V2.} A vision transformer cuts an image into patches, embeds each as a token and lets every token attend to every
other; the cost is quadratic in the number of tokens, which is prohibitive for dense prediction at high resolution. The Swin Transformer computes
self-attention only inside non-overlapping windows of $M\times M$ tokens and shifts the window grid by half a window between consecutive blocks, so
that information crosses window borders; patch-merging layers halve the resolution and double the channels between stages, which gives a
hierarchical pyramid like that of a convolutional network \citep{dosovitskiy2021vit,liu2021swin}. Swin Transformer V2 \citep{liu2022swinv2} adds three changes to scale capacity and
resolution: a residual post-norm combined with cosine attention to stabilise training, a log-spaced continuous position bias to transfer models
pre-trained at low resolution to high-resolution inputs, and SimMIM self-supervised pre-training. In cosine attention the similarity of a query
$q_i$ and a key $k_j$ is
\begin{equation}
\mathrm{Sim}(q_i,k_j)=\frac{\cos(q_i,k_j)}{\tau}+B_{ij},
\label{eq:bg_cosattn}
\end{equation}
with a learnable temperature $\tau$ and a relative position bias $B_{ij}$ produced by a small network from log-spaced relative coordinates.
Because attention is computed within windows of fixed size, an encoder pretrained on small images can be applied to a larger input with the same
window size; only the number of windows grows. The continuous bias additionally allows the window size itself to change.

\paragraph{ConvNeXt V2.} ConvNeXt is a convolutional network modernised with design choices from transformers: a patchify stem, large ($7\times7$)
depthwise convolutions, layer normalisation and inverted-bottleneck MLPs with GELU \citep{liu2022convnext}. ConvNeXt V2 \citep{woo2023convnextv2} co-designs the
architecture with masked-autoencoder pre-training: it proposes a fully convolutional masked autoencoder (FCMAE) and a Global Response
Normalization (GRN) layer that enhances inter-channel feature competition. For a feature map $X\in\mathrm{R}^{H\times W\times C}$, GRN computes
the spatial $\ell_2$ norm of each channel, normalises it by the mean over channels and rescales the features:
\begin{equation}
\begin{gathered}
G_i=\lVert X_{:,:,i}\rVert_2,\\
N_i=\frac{G_i}{\frac1C\sum_{j=1}^{C}G_j},\\
X_{:,:,i}\leftarrow\gamma_i\,X_{:,:,i}\,N_i+\beta_i+X_{:,:,i},
\end{gathered}
\label{eq:bg_grn}
\end{equation}
with learnable $\gamma$ and $\beta$. The encoder used here was pre-trained with FCMAE and then fine-tuned on ImageNet \citep{deng2009imagenet}
(Section~\ref{sec:method_v2}).

\paragraph{Nested decoder (UNet++).} A U-Net decoder upsamples the deepest features step by step and concatenates, at each scale, the encoder
features of the same resolution (skip connections). UNet++ \citep{zhou2018unetpp} replaces the plain skips by nested, dense skip pathways: node
$X^{i,j}$ at depth $i$ and position $j$ along the skip receives all previous nodes at the same depth and the upsampled node from the next depth,
\begin{equation}
X^{i,j}=\mathcal{H}\Big(\big[X^{i,0},X^{i,1},\dots,X^{i,j-1},\,\mathcal{U}(X^{i+1,j-1})\big]\Big),
\label{eq:bg_unetpp}
\end{equation}
where $\mathcal{H}$ is a convolution block, $\mathcal{U}$ upsampling and $[\cdot]$ concatenation. The intermediate nodes narrow the semantic gap
between encoder and decoder features before they are fused. UNet++ also proposes deep supervision of the intermediate outputs; the decoder here
uses only the final node. Because the transformer and ConvNeXt pyramids start at stride 4, the network of this study adds a short tail with two
shallow convolutional skips from the input to recover full resolution (Section~\ref{sec:methods}).

\paragraph{CBAM.} The Convolutional Block Attention Module \citep{woo2018cbam} reweights a feature map $F$ first by channel and then by
position:
\begin{equation}
\begin{gathered}
M_c(F)=\sigma\big(\mathrm{MLP}(\mathrm{AvgPool}(F))+\mathrm{MLP}(\mathrm{MaxPool}(F))\big),\\
M_s(F)=\sigma\big(f^{7\times7}([\mathrm{AvgPool}_c(F);\mathrm{MaxPool}_c(F)])\big),
\end{gathered}
\label{eq:bg_cbam}
\end{equation}
with $F'=M_c(F)\odot F$ and $F''=M_s(F')\odot F'$. Here $\sigma$ is the logistic function, the MLP is shared between the two pooled descriptors,
$\mathrm{AvgPool}_c$ and $\mathrm{MaxPool}_c$ pool over channels and $f^{7\times7}$ is a $7\times7$ convolution. In a suppressor the spatial map
can emphasise rib regions; the module is cheap and is placed after the decoder blocks.

\paragraph{AdamW, schedule and clipping.} Adam keeps exponential moving averages of the gradient $g_t$ and of its square, $m_t$ and $v_t$, and
scales each parameter's step by $1/\sqrt{\hat v_t}$. With an $L_2$ penalty added to the loss, the penalty gradient is rescaled in the same way, so
parameters with large gradient variance are regularised less. AdamW \citep{loshchilov2019adamw} decouples weight decay from the adaptive update:
\begin{equation}
\theta_t=\theta_{t-1}-\eta_t\Big(\frac{\hat m_t}{\sqrt{\hat v_t}+\epsilon}+\lambda\,\theta_{t-1}\Big),
\label{eq:bg_adamw}
\end{equation}
where $\hat m_t$ and $\hat v_t$ are the bias-corrected moments and $\lambda$ the decoupled weight-decay rate. The learning rate $\eta_t$ rises
linearly during a short warm-up of $t_w$ steps and then follows a cosine decay to zero at step $T$,
\begin{equation}
\begin{gathered}
\eta_t=\eta_{\max}\,\tfrac12\Big(1+\cos\Big(\pi\,\frac{t-t_w}{T-t_w}\Big)\Big),\\
t_w\le t\le T .
\end{gathered}
\label{eq:bg_cosine}
\end{equation}
A pretrained encoder is often given a smaller learning rate than a randomly initialised decoder, so that the pretrained features are adapted
rather than overwritten. Gradient clipping rescales the full gradient $g$ to $g\,\min(1,c/\lVert g\rVert_2)$ when its norm exceeds $c$, which
limits the damage of rare very large gradients.

\paragraph{EMA weights.} An exponential moving average of the weights,
\begin{equation}
\bar\theta_t=d_t\,\bar\theta_{t-1}+(1-d_t)\,\theta_t ,
\label{eq:bg_ema}
\end{equation}
is kept alongside the trained weights and used for validation and evaluation. Averaging over the last part of the trajectory smooths the noise of
stochastic optimisation and usually gives a slightly better and more stable model than the last iterate. The decay $d_t$ is small at the start and
grows to a fixed maximum, so that the average is not dominated by the random initialisation; normalisation statistics (buffers) are copied rather
than averaged.

\paragraph{Mixed precision.} Mixed-precision training runs most of the forward and backward pass in 16-bit floating point (fp16), which halves
memory traffic and uses the fast matrix units of current GPUs, while the master weights and the optimiser state stay in 32-bit precision. The
fp16 format has a small range: very small gradients underflow to zero and values above about $6.5\times10^4$ overflow. Dynamic loss scaling
multiplies the loss by a factor $S$ before the backward pass, divides the gradients by $S$ before clipping and the optimiser step, skips any step
whose gradients contain infinities or NaNs, and adapts $S$ during training. Loss scaling does not protect activations in the forward pass, and
fp16 overflow can still make a run diverge; Section~\ref{sec:experiments} reports a divergence attributed to fp16 overflow and how it was handled.

\section{Metrics}
\label{sec:bg_metrics}
Throughout this section $y$ is a reference image (for example a DES soft-tissue image), $\hat y$ an estimate of it, both scaled to $[0,1]$, $p$
indexes pixels and $R$ is a set of pixels (the whole image or a region).

\paragraph{Tone calibration.} A DES image and a radiograph differ in global tone (Section~\ref{sec:bg_des}), and a suppressor trained on other
data cannot know the tone of the reference. Fidelity metrics are therefore also computed after a per-image least-squares gain and offset,
\begin{equation}
\begin{gathered}
\hat y_{\mathrm{cal}}=\mathrm{clip}_{[0,1]}\big(a^\ast\hat y+c^\ast\big),\\
(a^\ast,c^\ast)=\arg\min_{a,c}\sum_p\big(a\,\hat y_p+c-y_p\big)^2 ,
\end{gathered}
\label{eq:bg_calib}
\end{equation}
which removes differences in brightness and contrast but not in structure. Section~\ref{sec:methods} calls these metrics ``calibrated''.

\paragraph{PSNR.} The mean squared error over $R$ and the peak signal-to-noise ratio are
\begin{equation}
\begin{gathered}
\mathrm{MSE}_R=\frac{1}{|R|}\sum_{p\in R}\big(\hat y_p-y_p\big)^2,\\
\mathrm{PSNR}_R=10\log_{10}\frac{L^2}{\mathrm{MSE}_R}\ \mathrm{dB},
\end{gathered}
\label{eq:bg_psnr}
\end{equation}
with dynamic range $L=1$. PSNR is a monotone transform of the MSE: it weights every pixel equally and is dominated by large errors, so it is
sensitive to the global tone and to a few badly wrong pixels. Restricting $R$ to the region where bone was removed (the bone region of
Section~\ref{sec:methods}) focuses the metric on the pixels that suppression should change.

\paragraph{SSIM.} The structural similarity index \citep{wang2004ssim} compares local means, variances and covariance of two images $u$ and $v$ in
a sliding window:
\begin{equation}
\begin{gathered}
\mathrm{SSIM}(u,v)=\frac{(2\mu_u\mu_v+C_1)\,(2\sigma_{uv}+C_2)}{(\mu_u^2+\mu_v^2+C_1)\,(\sigma_u^2+\sigma_v^2+C_2)},\\
C_1=(K_1L)^2,\ C_2=(K_2L)^2,
\end{gathered}
\label{eq:bg_ssim}
\end{equation}
with the customary $K_1=0.01$ and $K_2=0.03$. It is the product of a luminance term
$l=(2\mu_u\mu_v+C_1)/(\mu_u^2+\mu_v^2+C_1)$, a contrast term $c=(2\sigma_u\sigma_v+C_2)/(\sigma_u^2+\sigma_v^2+C_2)$ and a structure term
$s=(\sigma_{uv}+C_2/2)/(\sigma_u\sigma_v+C_2/2)$. The local values form an SSIM map, whose mean over $R$ is reported. SSIM ranges up to 1
(identical images); it is more sensitive to structure and local contrast than PSNR, and it depends on the window, which differs between
implementations (the original uses an $11\times11$ Gaussian window; the scikit-image function used here uses a $7\times7$ uniform window by
default).

\paragraph{MS-SSIM.} Multi-scale SSIM \citep{wang2003msssim} evaluates the contrast and structure terms at $M$ scales, obtained by repeated
low-pass filtering and down-sampling by two, and the luminance term only at the coarsest scale:
\begin{equation}
\mathrm{MS\text{-}SSIM}(u,v)=\big[l_M(u,v)\big]^{\alpha_M}\prod_{j=1}^{M}\big[c_j(u,v)\big]^{\beta_j}\big[s_j(u,v)\big]^{\gamma_j},
\label{eq:bg_msssim}
\end{equation}
with $M=5$ and the exponents fitted to human judgements in the original paper ($\beta_j=\gamma_j=0.0448$, $0.2856$, $0.3001$, $0.2363$, $0.1333$
and $\alpha_M=\beta_M$). It is less dominated by fine texture than single-scale SSIM.

\paragraph{Band-limited SSIM.} A DES soft-tissue image is filtered and noisier than the radiograph, so a suppressor that keeps the fine texture
of the radiograph can be penalised by full-resolution SSIM for texture that has nothing to do with bone. Band-limited SSIM compares the two images
after the same Gaussian low-pass filter $G_\sigma$,
\begin{equation}
\mathrm{SSIM}_\sigma(\hat y,y)=\mathrm{SSIM}\big(G_\sigma\ast\hat y,\;G_\sigma\ast y\big),
\label{eq:bg_blssim}
\end{equation}
which keeps the scales of ribs and soft-tissue structures and discards the finest texture. The values of $\sigma$ are fixed in the analysis plan
and applied identically to every method.

\paragraph{LPIPS.} The learned perceptual image patch similarity \citep{zhang2018lpips} measures distance in the feature space of a pretrained
network. With $\hat u^{l}_{hw}$ and $\hat v^{l}_{hw}$ the channel-normalised activations of layer $l$ at position $(h,w)$,
\begin{equation}
d(u,v)=\sum_l\frac{1}{H_lW_l}\sum_{h,w}\big\lVert w_l\odot\big(\hat u^{l}_{hw}-\hat v^{l}_{hw}\big)\big\rVert_2^2 ,
\label{eq:bg_lpips}
\end{equation}
where the channel weights $w_l$ were fitted to human similarity judgements. Lower is better. This study uses the AlexNet variant and replicates
the grey image to three channels. The network and the human judgements come from natural images, so LPIPS is a secondary, perceptual check
rather than a measure of diagnostic content.

\paragraph{Contrast-to-noise ratio of a box.} For a lesion with bounding box $\mathcal{B}$, let $\mathcal{R}$ be the ring obtained by enlarging
the box by half its width on the left and right and by half its height above and below, minus the box itself (clipped to the image). The
contrast-to-noise ratio is
\begin{equation}
\mathrm{CNR}=\frac{\bar x_{\mathcal{B}}-\bar x_{\mathcal{R}}}{\mathrm{sd}_{\mathcal{R}}(x)+\varepsilon},
\label{eq:bg_cnr}
\end{equation}
with means $\bar x$ and standard deviation $\mathrm{sd}$ over the indicated pixels and a small $\varepsilon$. The effect of suppression is the
change $\Delta\mathrm{CNR}=\mathrm{CNR}(\hat s)-\mathrm{CNR}(x)$ on the same box. Two effects are mixed in this number: removing rib edges from
the ring lowers $\mathrm{sd}_{\mathcal{R}}$ and raises the CNR, while moving part of the lesion into the bone image lowers the numerator. A
positive change therefore does not prove that the lesion was preserved, and the CNR is used as evidence that complements, but does not replace,
the insertion audit. A box also contains background around the lesion, so its mean is a diluted measure of lesion contrast.

\paragraph{Detection: FROC and CPM.} A detector outputs boxes with scores. A predicted box is matched to a ground-truth box if their
intersection over union, $\mathrm{IoU}=|\mathcal{B}_1\cap\mathcal{B}_2|/|\mathcal{B}_1\cup\mathcal{B}_2|$, reaches a threshold; predictions are
matched greedily in order of decreasing score and each ground-truth box at most once. For a score threshold $t$, let $\mathrm{TP}(t)$ and
$\mathrm{FP}(t)$ be the numbers of matched and unmatched predictions with score at least $t$, $N_{\mathrm{gt}}$ the number of ground-truth
lesions and $N_{\mathrm{img}}$ the number of images. The free-response ROC (FROC) curve plots the sensitivity
$\mathrm{TP}(t)/N_{\mathrm{gt}}$ against the mean number of false positives per image $\mathrm{FP}(t)/N_{\mathrm{img}}$ as $t$ varies. The
competition performance metric (CPM) averages the sensitivity at seven operating points,
\begin{equation}
\begin{gathered}
\mathrm{CPM}=\frac17\sum_{r\in\{\frac18,\frac14,\frac12,1,2,4,8\}}\mathrm{Sens}(r),\\
\mathrm{Sens}(r)=\max_{t:\ \mathrm{FP}(t)/N_{\mathrm{img}}\le r}\frac{\mathrm{TP}(t)}{N_{\mathrm{gt}}},
\end{gathered}
\label{eq:bg_cpm}
\end{equation}
as introduced in the ANODE09 study and named CPM by \citet{niemeijer2011cpm} \citep[see also][]{vanginneken2010anode09,setio2017luna16}. Tied
scores are treated as one threshold, so the curve does not depend on how a sort orders ties. FROC analysis suits lesion detection because an
image can contain several lesions and several false marks, which an image-level ROC curve cannot represent.

\paragraph{Average precision.} With precision $P(t)=\mathrm{TP}(t)/(\mathrm{TP}(t)+\mathrm{FP}(t))$ and recall
$\mathrm{Rec}(t)=\mathrm{TP}(t)/N_{\mathrm{gt}}$, the interpolated precision at recall $r$ is $P_{\mathrm{int}}(r)=\max_{r'\ge r}P(r')$ and the
average precision is the area under the interpolated precision--recall curve,
\begin{equation}
\mathrm{AP}=\sum_k\big(r_k-r_{k-1}\big)\,P_{\mathrm{int}}(r_k),
\label{eq:bg_ap}
\end{equation}
summed over the distinct recall levels $r_k$ (all-point interpolation). The mean over classes at an IoU threshold of 0.5 is reported as
mAP@50. Unlike the CPM, AP penalises false positives through precision rather than per image.

\paragraph{AUROC.} For a binary label with $m$ positive cases with scores $X_1,\dots,X_m$ and $n$ negative cases with scores $Y_1,\dots,Y_n$, the
area under the ROC curve equals the Mann--Whitney statistic
\begin{equation}
\begin{gathered}
\hat\theta=\frac{1}{mn}\sum_{i=1}^{m}\sum_{j=1}^{n}\psi(X_i,Y_j),\\
\psi(a,b)=\begin{cases}1 & a>b\\ \tfrac12 & a=b\\ 0 & a<b\end{cases},
\end{gathered}
\label{eq:bg_auc}
\end{equation}
the probability that a random positive case scores higher than a random negative one. It does not depend on a threshold or on the class
prevalence. For several labels (for example the findings of a multi-label classifier) the macro-AUROC is the unweighted mean of the per-label AUROCs.

\section{Uncertainty estimation}
\label{sec:bg_uncertainty}
A suppressor returns one bone image; a user also needs to know where that image may be wrong. Two kinds of uncertainty are usually distinguished:
\emph{aleatoric} uncertainty, which comes from ambiguity in the data and cannot be reduced by more training data (the rib-or-nodule ambiguity of
Section~\ref{sec:bg_cxr} is of this kind), and \emph{epistemic} uncertainty, which comes from limited data and model knowledge and grows on
inputs unlike the training data. This study uses a deep ensemble to obtain a per-pixel uncertainty score and split conformal prediction to turn
that score into intervals with a stated coverage.

\paragraph{Deep ensembles.} \citet{lakshminarayanan2017ensembles} train several networks independently from different random initialisations
(and data orders) and combine their predictions; they also use proper scoring rules and, optionally, adversarial training. For $K$ members
$f_1,\dots,f_K$, the ensemble prediction and spread of the bone image at pixel $p$ are
\begin{equation}
\begin{gathered}
\bar f_p(x)=\frac1K\sum_{k=1}^{K}f_{k,p}(x),\\
\sigma_p(x)=\Big(\frac{1}{K-1}\sum_{k=1}^{K}\big(f_{k,p}(x)-\bar f_p(x)\big)^2\Big)^{1/2}.
\end{gathered}
\label{eq:bg_ensemble}
\end{equation}
The mean is usually more accurate than a single member, and the spread is large where the members disagree, which happens mostly where the data
did not pin the function down. In this study each member outputs only a bone image, not a variance, so $\sigma$ reflects disagreement between
members rather than noise in the target. With few members $\sigma$ is itself a noisy estimate, and nothing guarantees that it is calibrated: a
spread of $0.01$ need not mean an error of about $0.01$. Whether the spread ranks pixels by their error is checked with the Spearman correlation
between spread and absolute error (Section~\ref{sec:bg_stats}) and with sparsification curves (Section~\ref{sec:method_v2}).

\paragraph{Split conformal prediction.} Conformal prediction turns any score into a prediction set with a guaranteed coverage
\citep{lei2018conformal,angelopoulos2023conformal}. In the split version the model is fixed and a calibration set
$(X_1,Y_1),\dots,(X_n,Y_n)$, disjoint from the training data, is used to compute nonconformity scores. For a regression output with an
uncertainty scale, a natural score is the normalised residual
\begin{equation}
S_i=\frac{|Y_i-\bar f(X_i)|}{\sigma(X_i)+\varepsilon},
\label{eq:bg_score}
\end{equation}
with a small $\varepsilon$ that keeps the interval open where the members agree exactly. Let $\hat q$ be the
$\lceil (n+1)(1-\alpha)\rceil$-th smallest of $S_1,\dots,S_n$. The interval for a new input is
\begin{equation}
C(X_{n+1})=\big[\bar f(X_{n+1})-\hat q\,(\sigma(X_{n+1})+\varepsilon),\;\bar f(X_{n+1})+\hat q\,(\sigma(X_{n+1})+\varepsilon)\big].
\label{eq:bg_conformal}
\end{equation}
For the bone image $Y$ is a bone pixel value; because $\hat s=x-\hat b$ and $x$ is known, the soft image has an interval of the same width.

\paragraph{Exchangeability and marginal coverage.} The guarantee rests on one assumption: the calibration pairs and the test pair
$(X_1,Y_1),\dots,(X_{n+1},Y_{n+1})$ are \emph{exchangeable}, that is, their joint distribution does not change when the indices are permuted.
Independent, identically distributed data are exchangeable. Under exchangeability the rank of $S_{n+1}$ among $S_1,\dots,S_{n+1}$ is uniformly
distributed, from which it follows that
\begin{equation}
\Pr\big(Y_{n+1}\in C(X_{n+1})\big)\ \ge\ 1-\alpha ,
\label{eq:bg_coverage}
\end{equation}
and, if the scores are almost surely distinct, the coverage is also at most $1-\alpha+1/(n+1)$ \citep{lei2018conformal}. Three limits of this
statement matter for this study. First, the coverage is \emph{marginal}: the probability is over the random calibration set and the random test
point together, so the interval covers $1-\alpha$ of test points on average, not $1-\alpha$ of the pixels of every image or of every region (for
example the pixels over a lesion). Second, when pixels are pooled as calibration units, the pixels of one image are dependent, so exchangeability
of pixels is an approximation; the resulting statement is about pixels drawn as in calibration, not about images (Section~\ref{sec:method_v2}).
Third, the guarantee holds only for test data from the calibration distribution. Calibrating on CT projections and testing on real radiographs
breaks exchangeability, so coverage on real images is measured, not guaranteed, and is reported as coverage under shift
(Chapter~\ref{ch:reliability}).

\section{Statistical methods}
\label{sec:bg_stats}
All comparisons in this thesis are \emph{paired}: every method is scored on the same units (images, lesion sites, patients), and inference is
based on the per-unit differences $d_i=a_i-b_i$ between two methods. Pairing removes the large variation between units that both methods share and
makes small differences detectable.

\paragraph{Percentile bootstrap.} The bootstrap \citep{efron1993bootstrap} estimates the sampling distribution of a statistic by resampling the
observed data. For the mean difference $\hat\theta=\frac1n\sum_i d_i$, draw $B$ samples of size $n$ with replacement from $d_1,\dots,d_n$,
compute the mean $\hat\theta^{\ast}_b$ of each, and take the empirical quantiles
\begin{equation}
\big[\hat\theta^{\ast}_{(\alpha/2)},\ \hat\theta^{\ast}_{(1-\alpha/2)}\big]
\label{eq:bg_boot}
\end{equation}
as a $1-\alpha$ confidence interval. The percentile interval needs no normality assumption; it assumes that the units are independent and that the
sample is large enough for the bootstrap distribution to approximate the sampling distribution.

\paragraph{Cluster bootstrap.} Units are often nested: several lesion sites per image, several images per patient. Sites in one image share the
image's anatomy and the model's behaviour on it, so they are positively correlated, and resampling sites as if independent gives intervals that
are too narrow. The cluster bootstrap resamples whole clusters $g=1,\dots,G$ with replacement and recomputes the statistic over all units in the
resampled clusters,
\begin{equation}
\hat\theta^{\ast}_b=\frac{\sum_{g\in\mathcal{S}_b}\sum_{i\in g}d_i}{\sum_{g\in\mathcal{S}_b}n_g},
\label{eq:bg_clusterboot}
\end{equation}
where $\mathcal{S}_b$ is the $b$-th resample of cluster indices (with repetitions) and $n_g$ the number of units in cluster $g$. The between-cluster
variation then enters the interval as it should. In this study, for example, lesion sites are clustered by image and images by patient
(Section~\ref{sec:experiments}).

\paragraph{Wilcoxon signed-rank test.} \citet{wilcoxon1945individual} proposed ranking the absolute paired differences. Zero differences are
dropped; the remaining $|d_i|$ are ranked from 1 to $n$ (ties receive average ranks), and the statistic is the sum of the ranks of the positive
differences,
\begin{equation}
\begin{gathered}
W^{+}=\sum_{i:\,d_i>0}\mathrm{rank}(|d_i|),\\
\mathrm{E}[W^{+}]=\frac{n(n+1)}{4},\\
\mathrm{Var}[W^{+}]=\frac{n(n+1)(2n+1)}{24}
\end{gathered}
\label{eq:bg_wilcoxon}
\end{equation}
under the null hypothesis that the differences are distributed symmetrically about zero (variance without ties). Large samples use the normal
approximation, small ones the exact distribution. The test is robust to outliers and does not assume normality, but it tests a shift in location
(the pseudo-median of the differences), not the mean; the bootstrap interval of the mean and the Wilcoxon $p$-value therefore answer slightly
different questions and are reported together.

\paragraph{Holm correction.} Testing $m$ hypotheses at level $\alpha$ each inflates the probability of at least one false rejection (the
family-wise error rate). The Holm procedure \citep{holm1979simple} orders the $p$-values $p_{(1)}\le\dots\le p_{(m)}$ and rejects
$H_{(1)},\dots,H_{(k-1)}$, where $k$ is the first index with $p_{(k)}>\alpha/(m-k+1)$. Equivalently, the adjusted $p$-values
\begin{equation}
\tilde p_{(i)}=\max_{j\le i}\ \min\big(1,\,(m-j+1)\,p_{(j)}\big)
\label{eq:bg_holm}
\end{equation}
are compared with $\alpha$. Holm controls the family-wise error rate under any dependence between the tests and rejects at least as much as
the Bonferroni correction. It is applied within each pre-declared family of hypotheses; confidence intervals are not adjusted.

\paragraph{One-sided non-inferiority.} A new method that adds a regulariser is not expected to be better than a reference on every metric; the
question is whether it is \emph{not worse by more than a tolerable amount}. For a metric where higher is better and a margin $\delta>0$ fixed in
advance, the hypotheses are
\begin{equation}
H_0:\ \Delta\le-\delta\qquad\text{against}\qquad H_1:\ \Delta>-\delta ,
\label{eq:bg_noninf}
\end{equation}
where $\Delta$ is the mean paired difference (new minus reference). $H_0$ is rejected at level $0.05$ if the lower limit of a one-sided 95\,\%
confidence interval for $\Delta$, which equals the lower limit of a two-sided 90\,\% interval, lies above $-\delta$. Three points are often
confused. Failing to show non-inferiority does not show inferiority; non-inferiority does not mean equality; and a margin chosen after seeing the
data makes the test meaningless, which is why the margins of this study were fixed in the analysis plan (Section~\ref{sec:experiments}).

\paragraph{DeLong test for paired AUCs.} When two classifiers score the same cases, their AUCs are correlated, and comparing them as if
independent wastes power. \citet{delong1988auc} use the structure of Equation~\eqref{eq:bg_auc}. For classifier $r$ define the placement values
\begin{equation}
\begin{gathered}
V^{r}_{10}(X_i)=\frac1n\sum_{j=1}^{n}\psi(X^{r}_i,Y^{r}_j),\\
V^{r}_{01}(Y_j)=\frac1m\sum_{i=1}^{m}\psi(X^{r}_i,Y^{r}_j),
\end{gathered}
\label{eq:bg_placements}
\end{equation}
whose means are both $\hat\theta^{r}$. With $S_{10}$ and $S_{01}$ the sample covariance matrices of the placement values over the positive and
the negative cases, the covariance of the AUC estimates is $S=S_{10}/m+S_{01}/n$, and
\begin{equation}
z=\frac{\hat\theta^{1}-\hat\theta^{2}}{\sqrt{S_{11}+S_{22}-2S_{12}}}
\label{eq:bg_delong}
\end{equation}
is compared with a standard normal distribution. The covariance term $S_{12}$ is what makes the test paired: the more the two classifiers agree
on which cases are hard, the smaller the variance of their difference.

\paragraph{Spearman correlation.} The Spearman coefficient is the Pearson correlation of the ranks of two variables; without ties it equals
\begin{equation}
\rho_S=1-\frac{6\sum_i D_i^2}{n(n^2-1)},
\label{eq:bg_spearman}
\end{equation}
where $D_i$ is the difference between the ranks of the two variables for unit $i$. It measures monotone rather than linear association and is
invariant to monotone transformations of either variable, which suits the comparison of an uncertainty score with an error, whose scales differ.

\paragraph{Training seeds.} The result of training a network depends on its random seed (initialisation, data order, augmentation). A difference
between two arms measured with one seed each confounds the effect of the arm with seed noise. Arms are therefore trained with several seeds,
seed results are averaged within each unit before a test, and the spread over seeds is reported alongside (Section~\ref{sec:methods}).

\paragraph{Pre-specified plans, development and confirmatory data.} A study that tries many models, metrics, subsets and tests and reports the
best has a much higher false-positive rate than its nominal $\alpha$, and a model chosen because it scored best on a data set will, on average,
score lower on new data from the same source (selection bias, or the winner's curse) \citep{simmons2011falsepositive,varoquaux2022failures}.
Two safeguards follow. First, development and confirmatory data are separated: every choice of architecture, weight, threshold and metric is
made on development data (here the validation projections and the SZCH development pairs), and the confirmatory data are scored only after
the choices are fixed. Second, the confirmatory analysis is pre-specified: hypotheses, primary metrics, tests, margins, hypothesis families and
the data lists are written down and frozen, with file hashes, before the confirmatory data are scored (Section~\ref{sec:experiments}).
Unlike a pre-registration, the plans and their hashes were kept in the project's version control and not deposited with a public registry,
so their timing rests on the commit history. A pre-specified plan does not forbid exploratory analyses; it makes the difference between
confirmatory and exploratory results visible. Any change
to the plan after freezing is a deviation and is reported as such (Appendix~\ref{app:deviations}), and a hypothesis that fails is reported as
failed.

\section{Related work}
\label{sec:related}
% PLAN (four subsections, about 1.5 pages). Describe a cited paper only as far as its abstract or title supports; check each claim against the PDF
% before submission.
%   2.1 Bone suppression: classical and CNN methods (suzuki2006mtann, yang2017cascade, gusarev2017deep); diffusion methods on real DES pairs
%       (chen2024bsdiff, sun2026bsldm, sun2025gllcm, sun2026debonedit). SZCH-X-Rays was introduced with BS-LDM and released with DeBoneDiT.
%   2.2 CT-supervised suppression (angaitkar2026anatomy; CT-RATE hamamci2026ctrate as the CT source of the projections we train on). What we reuse
%       (projections, released weights as comparator) and what we change (backbone, resolution, regulariser, evaluation).
%   2.3 Evaluation of restoration models for lesion fidelity: lesion insertion and removal, nodule generation (the NODE21 generation track,
%       sogancioglu2024node21), hallucination in generative restoration. [citations to be found, see paper/REFS_TODO.md]
%   2.4 Uncertainty and encoders: deep ensembles (lakshminarayanan2017ensembles), split conformal prediction (lei2018conformal,
%       angelopoulos2023conformal); modern encoders screened here (liu2022swinv2, woo2023convnextv2, simeoni2025dinov3, perezgarcia2025raddino).
\subsection{Bone suppression}
Early learning-based rib suppression used massive-training artificial neural networks \citep{suzuki2006mtann}; later work used convolutional
networks, for example a cascade of multi-scale networks in the gradient domain \citep{yang2017cascade} and other deep models
\citep{gusarev2017deep}. Recent methods train diffusion models on paired radiographs and DES soft-tissue images
\citep{chen2024bsdiff,sun2026bsldm,sun2025gllcm,sun2026debonedit}; the SZCH-X-Rays data set was compiled for BS-LDM \citep{sun2026bsldm} and made
public with DeBoneDiT \citep{sun2026debonedit}. BS-Diff trained on 120 pairs, augmented to 840, at $256\times256$ pixels; BS-LDM on 818 SZCH
pairs and 241 processed JSRT pairs at $1024\times1024$; GL-LCM and DeBoneDiT on the 741-pair SZCH version and JSRT at $1024\times1024$
(DeBoneDiT in a latent space of $4\times256\times256$).

\subsection{CT-supervised suppression}
\citet{angaitkar2026anatomy} build digitally reconstructed radiographs from chest CT, project bone, non-lung soft tissue and lung
separately, and train models that predict one component. Their projections are released on the Hugging Face Hub and are built from CT-RATE volumes
\citep{hamamci2026ctrate}: the release holds \num{21887} projections whose volume names are those of CT-RATE. We train on these projections and use the released weights as the main comparator. Our models differ in encoder and resolution (ConvNeXt-V2-T
at 512~px against SwinV2-B at 1024~px) and in sample budget, and add the twin regulariser, the ensemble and the audit.

\subsection{Evaluating lesion fidelity}
Learned image-to-image models can remove or invent content: distribution matching in translation networks can drop or hallucinate lesions
\citep{cohen2018distribution}, deep reconstruction can be unstable to small structures \citep{antun2020instabilities}, and hallucinations of
learned reconstruction have been defined and measured \citep{bhadra2021hallucinations,tivnan2024hallucination}. Inserting known features or signals into
images is an established way to test learned MRI reconstruction \citep{cheng2020addressing} and the effect of learned denoising on signal
detection \citep{li2021assessing}. For bone
suppression, simulated nodules were inserted to measure contrast after a hand-built clavicle suppression method \citep{hogeweg2013suppression},
and a phantom with a simulated nodule at many positions was used to measure contrast after bone suppression \citep{takagi2020dose}. These
studies measure lesion contrast; our audit measures the fraction of a known additive lesion signal that a learned suppressor moves into its
bone image, on real radiographs, and is shared by every model compared.

\subsection{Uncertainty and encoders}
Deep ensembles give a simple predictive-uncertainty estimate \citep{lakshminarayanan2017ensembles}, and split-conformal prediction turns an error score into
intervals with a coverage guarantee under exchangeability \citep{lei2018conformal,angelopoulos2023conformal}. We use both for the bone image
(Section~\ref{sec:experiments}). For the encoder we start from SwinV2-T \citep{liu2022swinv2} and screen ConvNeXt V2 \citep{woo2023convnextv2},
DINOv3 \citep{simeoni2025dinov3} and RAD-DINO \citep{perezgarcia2025raddino} under the same recipe: a modern convolutional encoder, a general
self-supervised vision transformer and a vision transformer pre-trained on chest radiographs.

